\documentclass[UTF8, 12pt, fontset=fandol, oneside]{ctexart}
\usepackage{researchnotes}

% 保留分册独立编译时的章号前缀。
\renewcommand{\thesection}{9.\arabic{section}}

\title{第九章\quad 数项级数}
\author{}
\date{}

\begin{document}

\maketitle

无穷级数的理论具有悠久的历史, 许多有趣的数学问题与它有关系, 在日常生活中我们也常常会遇到与它有关的问题. 如在古代人们为了计算圆的面积, 用圆内接正多边形的面积来逼近圆的面积, 该过程实际上就是通过一个无穷级数的部分和序列来求得圆面积具有任何精度的近似值. 另外, 无穷级数中有许多研究技巧与方法, 在许多学科中也有着广泛的用途. 当然, 若仅仅限于以上这些因素, 无穷级数未必能成为数学分析课程中专门讨论的一个对象. 我们在数学分析中研究无穷级数, 最主要的原因还是它对函数的研究起着重要的作用. 在本书的后续章节中可以看出这一点。

\section{数项级数的基本概念}

人们为了拓广函数类, 很原始的想法是将无穷多个函数相加起来. 我们注意到无穷多个函数相加时, 对定义域中固定的每一点来说, 也就是无穷多个数相加, 因此必须先研究形如
\begin{equation}
  a_1+\cdots+a_n+\cdots\triangleq \sum_{n=1}^{+\infty}a_n
  \tag{9.1.1}
\end{equation}
的和式. 另外, 在许多实际问题中, 我们也会遇到类似的和式, 如大家在中学学过的等比数列的和等. 在学习序列极限时, 我们也曾遇到过形如
\[
  x_n=1+\frac12+\frac13+\cdots+\frac1n,\quad n=1,2,\cdots,
\]
\[
  y_n=1+\frac1{2^2}+\frac1{3^2}+\cdots+\frac1{n^2},\quad n=1,2,\cdots
\]
的序列. 读者不难发现, 对 $n\geq1$, $x_n$ 其实就是和式 $\sum\limits_{k=1}^{+\infty}\dfrac1k$ 的前 $n$ 项之和; 而对 $n\geq1$, $y_n$ 则是 $\sum\limits_{k=1}^{+\infty}\dfrac1{k^2}$ 的前 $n$ 项之和. 通过前面的学习我们已经知道 $\lim\limits_{n\to+\infty}x_n=+\infty$, 而 $\lim\limits_{n\to+\infty}y_n$ 是存在的. 因此我们很自然地就说 $\sum\limits_{n=1}^{+\infty}\dfrac1n$ 是发散的, 而 $\sum\limits_{n=1}^{+\infty}\dfrac1{n^2}$ 是收敛的, 且 $\sum\limits_{n=1}^{+\infty}\dfrac1{n^2}=\lim\limits_{n\to+\infty}y_n$。

\subsection{数项级数的基本概念}

设 $a_1,a_2,\cdots,a_n,\cdots$ 为一个序列, 我们称
\[
  a_1+a_2+\cdots+a_n+\cdots
\]
为一个数项级数(简称级数), 记为 $\sum\limits_{n=1}^{+\infty}a_n$, 其中 $a_n(n=1,2,\cdots)$ 称为级数的通项, 并称 $S_n=\sum\limits_{k=1}^n a_k(n=1,2,\cdots)$ 为级数的前 $n$ 项部分和, 而称 $\{S_n\}$ 为部分和序列。

\noindent\textbf{定义 9.1.1}\quad 设 $\sum\limits_{n=1}^{+\infty}a_n$ 是一个数项级数. 如果它的部分和序列 $\{S_n\}$ 是收敛的, 即 $\lim\limits_{n\to\infty}S_n$ 存在, 则称级数 $\sum\limits_{n=1}^{+\infty}a_n$ 收敛, 并且记 $\sum\limits_{n=1}^{+\infty}a_n=\lim\limits_{n\to\infty}S_n$. 这时也称极限 $\lim\limits_{n\to\infty}S_n$ 为级数 $\sum\limits_{n=1}^{+\infty}a_n$ 的和. 如果 $\{S_n\}$ 是发散序列, 则称级数 $\sum\limits_{n=1}^{+\infty}a_n$ 发散. 特别地, 若 $\lim\limits_{n\to\infty}S_n=\infty(\pm\infty)$, 我们也称 $\sum\limits_{n=1}^{+\infty}a_n$ 发散到 $\infty(\pm\infty)$。

从以上定义知, 级数的敛散性本质上可以归结到序列的敛散性. 但我们以后会发现, 对级数的研究有许多特别的工具与方法, 似乎这些工具与方法不太容易直接从序列的研究中得到。

\noindent\textbf{例 9.1.1}\quad 对 $q\in\mathbb R$, 形如
\begin{equation}
  \sum_{n=1}^{+\infty}q^{n-1}=1+q+q^2+\cdots
  \tag{9.1.2}
\end{equation}
的级数称为等比级数(或几何级数). 试讨论该级数的敛散性。

\noindent\textbf{解}\quad 对于 $\forall n\in\mathbb N$, 该级数的前 $n$ 项部分和为
\[
  S_n(q)=\sum_{k=1}^n q^{k-1}=\frac{1-q^n}{1-q}\quad(q\neq1).
\]
当 $|q|<1$ 时, 我们有 $\lim\limits_{n\to\infty}S_n(q)=\dfrac1{1-q}$. 因此此时级数 (9.1.2) 收敛, 并且有
\[
  \sum_{n=1}^{+\infty}q^{n-1}=\frac1{1-q}.
\]
当 $q=-1$ 时, $S_n(-1)=\dfrac{1-(-1)^n}{2}$, 显然 $\lim\limits_{n\to\infty}S_n(q)$ 不存在. 因此此时级数 (9.1.2) 发散, 即 $\sum\limits_{n=1}^{+\infty}(-1)^{n-1}$ 发散。

当 $q=1$ 时, 有 $S_n(1)=n(n=1,2,\cdots)$, 因此 $\lim\limits_{n\to\infty}S_n(1)=+\infty$, 从而此时级数 (9.1.2) 发散。

当 $|q|>1$ 时, $\lim\limits_{n\to\infty}S_n(q)=\lim\limits_{n\to\infty}\dfrac{q^n-1}{q-1}=\infty$, 因此此时级数 (9.1.2) 发散。

所以当 $|q|<1$ 时, 等比级数 (9.1.2) 收敛; 当 $|q|\geq1$ 时, 级数 (9.1.2) 发散. 特别地, 当 $q<-1$ 时, 级数 (9.1.2) 发散到 $\infty$; 当 $q>1$ 时, 级数 (9.1.2) 发散到 $+\infty$。

\noindent\textbf{例 9.1.2}\quad 对固定的 $n_0\in\mathbb N$, 证明数项级数 $\sum\limits_{n=1}^{+\infty}\dfrac1{n(n+n_0)}$ 收敛, 并求其值。

\noindent\textbf{解}\quad 对于 $\forall n\in\mathbb N$, 该级数的前 $n$ 项部分和 $S_n$ 可以表示成
\[
\begin{aligned}
  S_n
  &=\sum_{k=1}^n\frac1{k(k+n_0)}
    =\sum_{k=1}^n\frac1{n_0}\left(\frac1k-\frac1{k+n_0}\right)\\
  &=\frac1{n_0}\left(\sum_{k=1}^{n_0}\frac1k+\sum_{k=n_0+1}^n\frac1k-\sum_{k=n_0+1}^{n+n_0}\frac1k\right)\\
  &=\frac1{n_0}\sum_{k=1}^{n_0}\frac1k+\frac1{n_0}\sum_{k=n+1}^{n+n_0}\frac1k.
\end{aligned}
\]
由于
\[
  0<\frac1{n_0}\sum_{k=n+1}^{n+n_0}\frac1k\leq \frac1{n_0}\cdot\frac{n_0}{n+1}\to0\quad(n\to\infty),
\]
因此 $\sum\limits_{n=1}^{+\infty}\dfrac1{n(n+n_0)}$ 收敛, 并且
\[
  \sum_{n=1}^{+\infty}\frac1{n(n+n_0)}
  =\lim_{n\to\infty}S_n=\frac1{n_0}\sum_{k=1}^{n_0}\frac1k.
\]

\noindent\textbf{例 9.1.3}\quad 证明 $\sum\limits_{n=0}^{+\infty}\dfrac1{n!}=e$, 并证明 $e$ 是无理数。

\noindent\textbf{证明}\quad 记 $S_n=\sum\limits_{k=0}^{n-1}\dfrac1{k!}(n=1,2,\cdots)$. 对任意固定的 $n>1$, 任取 $k>n$, 有
\[
  e>\left(1+\frac1k\right)^k
  >1+1+\frac1{2!}\left(1-\frac1k\right)+\cdots+\frac1{(n-1)!}\left(1-\frac1k\right)\cdots\left(1-\frac{n-2}{k}\right).
\]
在上式中令 $k\to\infty$, 得
\[
  1+1+\frac1{2!}+\frac1{3!}+\cdots+\frac1{(n-1)!}\leq e.
\]
再由
\[
  \left(1+\frac1{n-1}\right)^{n-1}<1+1+\frac1{2!}+\frac1{3!}+\cdots+\frac1{(n-1)!}\leq e,
\]
得 $\lim\limits_{n\to\infty}S_n=e$. 由级数收敛的定义知 $e=\sum\limits_{n=0}^{+\infty}\dfrac1{n!}$。

对任意正整数 $n,m$, 有
\[
\begin{aligned}
  0<S_{n+m}-S_n
  &=\frac1{n!}+\frac1{(n+1)!}+\cdots+\frac1{(n+m-1)!}\\
  &\leq\frac1{n!}\left[1+\frac1{n+1}+\frac1{(n+1)^2}+\cdots+\frac1{(n+1)^{m-1}}\right]\\
  &<\frac1{n!}\cdot\frac{n+1}{n}.
\end{aligned}
\]
在上式中固定 $n$, 让 $m\to+\infty$, 得
\begin{equation}
  0<e-S_n\leq\frac1{n!}\cdot\frac{n+1}{n}.
  \tag{9.1.3}
\end{equation}
从上式我们容易推出 $e$ 是无理数. 事实上, 若 $e$ 是有理数, 则存在既约正整数 $p,q$, 使得 $e=\dfrac qp$. 取定 $n\geq\max\{2,p\}+1$, 在 (9.1.3) 后一不等式的两边乘上 $(n-1)!$, 则其左边是一个正整数, 而右边是一个真分数. 此矛盾便证明了 $e$ 是一个无理数。

在第一册里, 我们曾利用指数函数 $e^x$ 的带拉格朗日余项的泰勒公式证明过 $e$ 是无理数. 在这里, 我们仅仅利用数项级数收敛的定义就证明了该结论。

关于数项级数我们有以下简单的性质, 其证明作为练习留给读者。

\noindent\textbf{性质 9.1.1}\quad 改变数项级数有限项的值后得到的新级数与原级数敛散性相同。

\noindent\textbf{性质 9.1.2}\quad 对于任意常数 $k\neq0$, 数项级数 $\sum\limits_{n=1}^{+\infty}a_n$ 与 $\sum\limits_{n=1}^{+\infty}ka_n$ 的敛散性相同。

\noindent\textbf{性质 9.1.3}\quad 设 $\alpha,\beta\in\mathbb R$, 数项级数 $\sum\limits_{n=1}^{+\infty}a_n$ 和 $\sum\limits_{n=1}^{+\infty}b_n$ 都收敛, 则数项级数 $\sum\limits_{n=1}^{+\infty}(\alpha a_n+\beta b_n)$ 也收敛, 并且成立下述等式:
\[
  \sum_{n=1}^{+\infty}(\alpha a_n+\beta b_n)
  =\alpha\sum_{n=1}^{+\infty}a_n+\beta\sum_{n=1}^{+\infty}b_n.
\]

在性质 9.1.3 中, 特别地, 当 $\beta=0$ 时, 对于收敛的数项级数 $\sum\limits_{n=1}^{+\infty}a_n$, 有
\[
  \sum_{n=1}^{+\infty}\alpha a_n=\alpha\sum_{n=1}^{+\infty}a_n\quad(\alpha\in\mathbb R).
\]

\subsection{柯西准则}

对于很多数项级数, 并不能简单地计算出它的部分和, 因此很难从定义来判别其敛散性, 从而我们必须直接从数项级数的本身来研究它是否收敛. 在下面的两节中我们将详细地研究这个问题. 在这里, 我们根据数项级数敛散性的定义, 给出以下判别级数收敛的柯西准则。

\noindent\textbf{定理 9.1.4 (柯西准则)}\quad 设 $\sum\limits_{n=1}^{+\infty}a_n$ 是一个数项级数, 则它收敛的充分必要条件是: 对于 $\forall\varepsilon>0$, $\exists N>0$, 当 $n>m>N$ 时, 有
\[
  \left|\sum_{k=m+1}^n a_k\right|
  =|a_{m+1}+a_{m+2}+\cdots+a_n|<\varepsilon.
\]
此定理的证明可以从数项级数收敛的定义直接推出, 故省略。

若数项级数 $\sum\limits_{n=1}^{+\infty}a_n$ 收敛, 由柯西准则, 对于 $\forall\varepsilon>0$, $\exists N>0$, 特别地取 $n,n-1\geq N$ 时, 我们有 $|a_n|<\varepsilon$, 从而我们有下面的重要结论。

\noindent\textbf{推论}\quad 设 $\sum\limits_{n=1}^{+\infty}a_n$ 是一个数项级数, 若它收敛, 则必有
\[
  \lim_{n\to\infty}a_n=0.
\]

显然, 推论的逆命题是不成立的, 如级数 $\sum\limits_{n=1}^{+\infty}\dfrac1n$(称为调和级数)是发散的, 但是有 $\lim\limits_{n\to\infty}\dfrac1n=0$。

\noindent\textbf{例 9.1.4}\quad 证明数项级数 $\sum\limits_{n=1}^{+\infty}n\ln\left(1+\dfrac1n\right)$ 发散。

\noindent\textbf{证明}\quad 由于
\[
  \lim_{n\to\infty}n\ln\left(1+\frac1n\right)
  =\lim_{n\to\infty}\ln\left(1+\frac1n\right)^n=1,
\]
因此由定理 9.1.4 的推论即知, 该数项级数是发散的。

\noindent\textbf{例 9.1.5}\quad 证明: 当 $\alpha\neq k\pi(k\in\mathbb Z)$ 时, 数项级数 $\sum\limits_{n=1}^{+\infty}\sin n\alpha$ 发散。

\noindent\textbf{证明}\quad 倘若 $\sum\limits_{n=1}^{+\infty}\sin n\alpha$ 收敛, 由推论 9.1.5 知 $\lim\limits_{n\to\infty}\sin n\alpha=0$. 由于对 $\forall n\in\mathbb N$, 我们有
\[
  \sin(n+1)\alpha=\cos\alpha\sin n\alpha+\cos n\alpha\sin\alpha,
\]
因此我们推出 $\lim\limits_{n\to\infty}\cos n\alpha=0$. 显然这与恒等式 $\sin^2n\alpha+\cos^2n\alpha=1$ 矛盾. 因此数项级数 $\sum\limits_{n=1}^{+\infty}\sin n\alpha$ 发散。

\section{正项级数}

\subsection{比较判别法}

对于数项级数 $\sum\limits_{n=1}^{+\infty}a_n$ 而言, 与无穷积分一样, 也有绝对收敛和条件收敛的概念。

\noindent\textbf{定义 9.2.1}\quad 设 $\sum\limits_{n=1}^{+\infty}a_n$ 为一个数项级数. 如果 $\sum\limits_{n=1}^{+\infty}|a_n|$ 收敛, 则称 $\sum\limits_{n=1}^{+\infty}a_n$ 绝对收敛. 如果 $\sum\limits_{n=1}^{+\infty}a_n$ 收敛, 但 $\sum\limits_{n=1}^{+\infty}|a_n|$ 发散, 则称 $\sum\limits_{n=1}^{+\infty}a_n$ 条件收敛。

由柯西准则及三角不等式易于推出: 若 $\sum\limits_{n=1}^{+\infty}a_n$ 绝对收敛, 则 $\sum\limits_{n=1}^{+\infty}a_n$ 必定收敛。

由定义 9.2.1 可看出, 一个级数是否为绝对收敛的问题可以转化为所谓的正项级数的敛散性问题。

\noindent\textbf{定义 9.2.2}\quad 设 $\sum\limits_{n=1}^{+\infty}a_n$ 是一个数项级数. 若对于 $\forall n\in\mathbb N$, 有 $a_n\geq0$, 则称 $\sum\limits_{n=1}^{+\infty}a_n$ 为一个正项级数。

对于正项级数, 以下结论是显然的。

\noindent\textbf{定理 9.2.1}\quad 设 $\sum\limits_{n=1}^{+\infty}a_n$ 为一个正项级数, 则它收敛的充分必要条件是它的部分和序列是有界的. 如果 $\sum\limits_{n=1}^{+\infty}a_n$ 发散, 则它必发散到 $+\infty$。

判别正项级数敛散性最基本的方法是如下的比较判别法。

\noindent\textbf{定理 9.2.2(比较判别法)}\quad 设 $\sum\limits_{n=1}^{+\infty}a_n$ 与 $\sum\limits_{n=1}^{+\infty}b_n$ 为两个正项级数, $c_1,c_2$ 是两个正数. 若存在 $N\in\mathbb N$, 当 $n>N$ 时, 有
\begin{equation}
  c_1a_n\leq c_2b_n,
  \tag{9.2.1}
\end{equation}
则

(1) 当 $\sum\limits_{n=1}^{+\infty}b_n$ 收敛时, $\sum\limits_{n=1}^{+\infty}a_n$ 也收敛;

(2) 当 $\sum\limits_{n=1}^{+\infty}a_n$ 发散时, $\sum\limits_{n=1}^{+\infty}b_n$ 也发散。

\noindent\textbf{证明}\quad 因为一个级数改变有限多项的值后得到的级数与原级数具有相同的敛散性, 因此不妨设 (9.2.1) 式对一切 $n\geq1$ 都成立。

对于 $n\in\mathbb N$, 分别记 $S_n$ 和 $S'_n$ 为 $\sum\limits_{n=1}^{+\infty}a_n$ 和 $\sum\limits_{n=1}^{+\infty}b_n$ 的前 $n$ 项部分和, 则由 (9.2.1) 式有
\[
  c_1S_n\leq c_2S'_n\quad(n=1,2,\cdots).
\]
(1) 当 $\sum\limits_{n=1}^{+\infty}b_n$ 收敛时, 由定理 9.2.1 知, $\exists M>0$, 使得对于 $\forall n\in\mathbb N$, 有
\[
  S'_n\leq M.
\]
因此, 对一切 $n\in\mathbb N$, 有
\[
  S_n\leq\frac{c_2}{c_1}M.
\]
由定理 9.2.1 知, $\sum\limits_{n=1}^{+\infty}a_n$ 也收敛。

(2) 当 $\sum\limits_{n=1}^{+\infty}a_n$ 发散时, 有 $S_n\to+\infty$, 从而有
\[
  \lim_{n\to\infty}S'_n\geq\lim_{n\to\infty}\frac{c_1}{c_2}S_n=+\infty.
\]
因此 $\sum\limits_{n=1}^{+\infty}b_n$ 也发散。

由极限的保号性及定理 9.2.2, 我们有以下的推论。

\noindent\textbf{推论}\quad 设正项级数 $\sum\limits_{n=1}^{+\infty}a_n$ 和 $\sum\limits_{n=1}^{+\infty}b_n(b_n\neq0,\forall n\in\mathbb N)$ 满足
\[
  \lim_{n\to\infty}\frac{a_n}{b_n}=l,
\]
则

(1) 当 $0<l<+\infty$ 时, $\sum\limits_{n=1}^{+\infty}a_n$ 和 $\sum\limits_{n=1}^{+\infty}b_n$ 具有相同的敛散性;

(2) 当 $l=0$ 时, 若 $\sum\limits_{n=1}^{+\infty}b_n$ 收敛, 则 $\sum\limits_{n=1}^{+\infty}a_n$ 也收敛;

(3) 当 $l=+\infty$ 时, 若 $\sum\limits_{n=1}^{+\infty}b_n$ 发散, 则 $\sum\limits_{n=1}^{+\infty}a_n$ 也发散。

\noindent\textbf{例 9.2.1}\quad 证明: 正项级数 $\sum\limits_{n=1}^{+\infty}\dfrac1{n^p}$ 当 $p\leq1$ 时发散, 当 $p>1$ 时收敛。

\noindent\textbf{证明}\quad 当 $p=1$ 时, 我们已经知道 $\sum\limits_{n=1}^{+\infty}\dfrac1n$ 是发散的. 当 $p<1$ 时, 对于 $\forall n>1$, 有 $\dfrac1n<\dfrac1{n^p}$, 因此 $\sum\limits_{n=1}^{+\infty}\dfrac1{n^p}$ 发散。

为了证明 $\sum\limits_{n=1}^{+\infty}\dfrac1{n^p}$ 当 $p>1$ 时收敛, 我们首先注意到以下简单事实: 因为一个正项级数 $\sum\limits_{n=1}^{+\infty}a_n$ 的部分和序列 $\{S_n\}$ 是单调上升的, 因此若它有一个子列是有界的, 则 $\sum\limits_{n=1}^{+\infty}a_n$ 必收敛。

注意到 $\left\{\dfrac1{n^p}\right\}$ 是一个单调下降序列, 对于 $\forall k\in\mathbb N$, 我们有
\[
\begin{aligned}
  S_{2^{k+1}-1}
  &=1+\left(\frac1{2^p}+\frac1{3^p}\right)
    +\left(\frac1{4^p}+\cdots+\frac1{7^p}\right)+\cdots\\
  &\quad+\left[\frac1{(2^k)^p}+\frac1{(2^k+1)^p}+\cdots+\frac1{(2^{k+1}-1)^p}\right]\\
  &\leq1+2\cdot\frac1{2^p}+2^2\cdot\frac1{2^{2p}}+\cdots+2^k\cdot\frac1{2^{kp}}\\
  &=\sum_{j=0}^k2^{(1-p)j}<\frac1{1-2^{1-p}}.
\end{aligned}
\]
上式说明, 当 $p>1$ 时, $\sum\limits_{n=1}^{+\infty}\dfrac1{n^p}$ 的部分和序列 $\{S_n\}$ 的子列 $\{S_{2^{k+1}-1}\}$ 有界, 因此该级数是收敛的。

在正项级数敛散性的讨论中, 几何级数 $\sum\limits_{n=1}^{+\infty}q^n$ 与 $\sum\limits_{n=1}^{+\infty}\dfrac1{n^p}$ 是两个常用的参考级数. 不仅许多其他的级数常用它们来作比较, 而且很多判别法也是建立在这两个级数的基础上的。

\noindent\textbf{例 9.2.2}\quad 设 $F_1=1,F_2=2,F_n=F_{n-1}+F_{n-2}(n=3,4,\cdots)$, 证明正项级数 $\sum\limits_{n=1}^{+\infty}\dfrac1{F_n}$ 收敛。

\noindent\textbf{证明}\quad 对于 $n=1,2,\cdots$, 记 $S_n=\sum\limits_{k=1}^n\dfrac1{F_k}$. 对任意 $n\geq3$, 我们有
\[
  F_{n-2}<F_{n-1}<2F_{n-2},
\]
因此有
\[
  F_n>F_{n-1}+\frac12F_{n-1}=\frac32F_{n-1}.
\]
利用上式容易看出
\[
  \lim_{n\to\infty}F_n=+\infty.
\]
同时我们还可以推出
\[
  \sum_{k=3}^n\frac1{F_k}\leq\frac23\sum_{k=3}^n\frac1{F_{k-1}},
\]
即
\[
  S_n-1-\frac12<\frac23S_n-\frac23-\frac23F_n^{-1}.
\]
将上式化简得
\[
  S_n<\frac52-2F_n^{-1}<3.
\]
因此 $\sum\limits_{n=1}^{+\infty}\dfrac1{F_n}$ 收敛。

\noindent\textbf{例 9.2.3}\quad 试讨论数项级数 $\sum\limits_{n=3}^{+\infty}\ln\cos\dfrac{\pi}{n}$ 的敛散性。

\noindent\textbf{解}\quad 记 $a_n=\ln\cos\dfrac{\pi}{n}$, 则 $a_n<0$. 由
\[
  \ln\cos\frac{\pi}{n}
  =\ln\left(1+\cos\frac{\pi}{n}-1\right)
  \sim \cos\frac{\pi}{n}-1\sim-\frac{\pi^2}{2n^2}\quad(n\to\infty)
\]
及 $\sum\limits_{n=1}^{+\infty}\dfrac{\pi^2}{2n^2}$ 收敛, 可知正项级数 $\sum\limits_{n=1}^{+\infty}\left(-\ln\cos\dfrac{\pi}{n}\right)$ 收敛, 因此原级数收敛。

\noindent\textbf{例 9.2.4}\quad 试讨论正项级数 $\sum\limits_{n=1}^{+\infty}\left[1-\left(\dfrac{n-1}{n+1}\right)^k\right]^p(k>0,p>0)$ 的敛散性。

\noindent\textbf{解}\quad 记 $a_n=\left[1-\left(\dfrac{n-1}{n+1}\right)^k\right]^p(n=1,2,\cdots)$. 由
\[
\begin{aligned}
  \left(\frac{n-1}{n+1}\right)^k
  &=\left(1-\frac1n\right)^k\left(1+\frac1n\right)^{-k}\\
  &=\left[1-\frac{k}{n}+o\left(\frac1n\right)\right]\left[1-\frac{k}{n}+o\left(\frac1n\right)\right]\\
  &=1-\frac{2k}{n}+o\left(\frac1n\right)\quad(n\to\infty),
\end{aligned}
\]
得
\[
  a_n=\left[\frac{2k}{n}+o\left(\frac1n\right)\right]^p\sim\frac{(2k)^p}{n^p}\quad(n\to\infty).
\]
由比较判别法知, 当 $p>1$ 时, $\sum\limits_{n=1}^{+\infty}a_n$ 收敛; 当 $p\leq1$ 时, $\sum\limits_{n=1}^{+\infty}a_n$ 发散。

\noindent\textbf{例 9.2.5}\quad 讨论正项级数 $\sum\limits_{n=1}^{+\infty}\dfrac{n^{n-2}}{e^n n!}$ 的敛散性。

\noindent\textbf{解}\quad 记 $a_n=\dfrac{n^{n-2}}{e^n n!}(n=1,2,\cdots)$. 对任意 $n\geq2$, 我们有
\[
  \frac{a_{n+1}}{a_n}
  =\frac{\left(1+\frac1n\right)^{n-2}}{e}
  <\frac{\left(1+\frac1n\right)^{n-2}}{\left(1+\frac1n\right)^n}
  =\frac{\dfrac1{(n+1)^2}}{\dfrac1{n^2}}
 =\frac{b_{n+1}}{b_n},
\]
其中 $b_n=\dfrac1{n^2}$。

由上面的不等式我们推出
\[
\begin{aligned}
  \frac{a_n}{a_2}
  &=\frac{a_n}{a_{n-1}}\frac{a_{n-1}}{a_{n-2}}\cdots\frac{a_3}{a_2}\\
  &<\frac{b_n}{b_{n-1}}\frac{b_{n-1}}{b_{n-2}}\cdots\frac{b_3}{b_2}
  =\frac{b_n}{b_2}=\frac4{n^2}.
\end{aligned}
\]
因此当 $n\geq2$ 时, 有 $a_n<\dfrac2{e^2}\cdot\dfrac1{n^2}$. 由定理 9.2.2 知 $\sum\limits_{n=1}^{+\infty}\dfrac{n^{n-2}}{e^n n!}$ 收敛。

\subsection{达朗贝尔判别法与柯西判别法}

以下两个判别法给出了从正项级数 $\sum a_n$ 本身的性质来判断其敛散性的方法, 但这些判别法的本质还是基于将正项级数与几何级数进行比较而得到的。

\noindent\textbf{定理 9.2.3 (达朗贝尔(D'Alembert)判别法)}\quad 设 $\sum a_n(a_n\neq0,n=1,2,\cdots)$ 为正项级数, 则

(1) 当
\[
  \overline{\lim}_{n\to\infty}\frac{a_{n+1}}{a_n}=\bar r<1
\]
时, $\sum a_n$ 收敛;

(2) 当
\[
  \underline{\lim}_{n\to\infty}\frac{a_{n+1}}{a_n}=\underline r>1
\]
时, $\sum a_n$ 发散。

\noindent\textbf{证明}\quad 我们先证 (1). 取 $r_1$, 使得其满足 $\bar r<r_1<1$, 则存在 $N_1\in\mathbb{N}$, 当 $n>N_1$ 时, 有 $\dfrac{a_n}{a_{n-1}}<r_1$. 因此有
\[
  a_n=\frac{a_n}{a_{n-1}}\frac{a_{n-1}}{a_{n-2}}\cdots\frac{a_{N_1+1}}{a_{N_1}}a_{N_1}
  <r_1^{n-N_1}a_{N_1}=C_1r_1^n,
\]
其中 $C_1=a_{N_1}r_1^{-N_1}$. 由于 $\sum r_1^n$ 是收敛的, 因此由定理 9.2.2 知 $\sum a_n$ 也收敛。

现证 (2). 取 $r_2$, 使得其满足 $\underline r>r_2>1$, 则存在 $N_2\in\mathbb{N}$, 当 $n>N_2$ 时, 有 $\dfrac{a_n}{a_{n-1}}>r_2$. 因此有
\[
  a_n=\frac{a_n}{a_{n-1}}\frac{a_{n-1}}{a_{n-2}}\cdots\frac{a_{N_2+1}}{a_{N_2}}a_{N_2}
  >r_2^{n-N_2}a_{N_2}=C_2r_2^n,
\]
其中 $C_2=a_{N_2}r_2^{-N_2}$. 由此推出 $\lim\limits_{n\to\infty}a_n=+\infty$, 因此 $\sum a_n$ 发散. 证毕。

对于 $\forall p>0$, 我们有
\[
  \lim_{n\to\infty}\frac{1/(n+1)^p}{1/n^p}
  =\lim_{n\to\infty}\left(\frac{n}{n+1}\right)^p=1,
\]
因此用达朗贝尔判别法我们无法判别正项级数 $\sum\dfrac1{n^p}$ 的敛散性. 换句话说, 当定理 9.2.3 中的 $\underline r$ 和 $\bar r$ 等于 1 时, 达朗贝尔判别法是失效的。

\noindent\textbf{例 9.2.6}\quad 试讨论正项级数 $\sum\limits_{n=1}^{+\infty}\dfrac{x^n n!}{n^n}$ 的敛散性, 其中 $x\geq0$。

\noindent\textbf{解}\quad 对于 $\forall n\in\mathbb{N}$, 记 $a_n=\dfrac{x^n n!}{n^n}$, 则有
\[
  \frac{a_{n+1}}{a_n}
  =\frac{\dfrac{x^{n+1}(n+1)!}{(n+1)^{n+1}}}{\dfrac{x^n n!}{n^n}}
  =\frac{x(n+1)}{\left(1+\frac1n\right)^n(n+1)}
  =\frac{x}{\left(1+\frac1n\right)^n}\to\frac{x}{e}\quad(n\to\infty).
\]
因此由达朗贝尔判别法知, $\sum\limits_{n=1}^{+\infty}\dfrac{x^n n!}{n^n}$ 当 $0\leq x<e$ 时收敛, 当 $x>e$ 时发散。

当 $x=e$ 时, 由于
\[
  (n+1)^n<e^n n!\quad(n>1),
\]
从而 $\lim\limits_{n\to\infty}a_n\neq0$, 因此该级数发散。

\noindent\textbf{定理 9.2.4 (柯西判别法)}\quad 设 $\sum a_n$ 为一个正项级数, 记
\[
  \overline{\lim}_{n\to\infty}\sqrt[n]{a_n}=r,
\]
则

(1) $r<1$ 时, $\sum a_n$ 收敛;

(2) $r>1$ 时, $\sum a_n$ 发散。

\noindent\textbf{证明}\quad (1) 当 $r<1$ 时, 取 $r_1$ 满足 $r<r_1<1$, 则 $\exists N_1$, 当 $n>N_1$ 时, 有 $\sqrt[n]{a_n}<r_1$, 即 $a_n<r_1^n$. 而 $\sum r_1^n$ 收敛, 因此 $\sum a_n$ 收敛。

(2) 当 $r>1$ 时, 取 $r_2$ 满足 $r>r_2>1$, 则对于 $\forall N\in\mathbb{N}$, $\exists n'>N$, 使得 $\sqrt[n']{a_{n'}}>r_2$, 即 $a_{n'}>r_2^{n'}$. 这说明 $\lim\limits_{n\to\infty}a_n\neq0$, 因此 $\sum a_n$ 发散. 证毕。

对于 $\forall p>0$, 我们有 $\lim\limits_{n\to\infty}\sqrt[n]{\dfrac1{n^p}}=1$, 因此用柯西判别法我们也无法判别正项级数 $\sum\dfrac1{n^p}$ 的敛散性. 换句话说, 当定理 9.2.4 中的 $r$ 为 1 时, 柯西判别法也是失效的。

从理论上来说, 凡是能用达朗贝尔判别法来判别敛散性的正项级数均可由柯西判别法来加以判别. 事实上, 对一般正项级数 $\sum a_n(a_n>0(n=1,2,\cdots))$, 下面的不等式总是成立的:
\[
  \underline{\lim}_{n\to\infty}\frac{a_{n+1}}{a_n}
  \leq \underline{\lim}_{n\to\infty}\sqrt[n]{a_n}
  \leq \overline{\lim}_{n\to\infty}\sqrt[n]{a_n}
  \leq \overline{\lim}_{n\to\infty}\frac{a_{n+1}}{a_n}.
\]
上述不等式的证明留给读者。

\noindent\textbf{例 9.2.7}\quad 试讨论正项级数 $\sum\limits_{n=1}^{+\infty}\left(\sqrt[n]{n}-\sin\dfrac1n\right)^{n^2}$ 的敛散性。

\noindent\textbf{解}\quad 记 $a_n=\left(\sqrt[n]{n}-\sin\dfrac1n\right)^{n^2}$. 由
\[
\begin{aligned}
  \sqrt[n]{n}-\sin\frac1n
  &=e^{\frac{\ln n}{n}}-\sin\frac1n\\
  &=1+\frac{\ln n}{n}+o\left(\frac{\ln n}{n}\right)-\left[\frac1n+o\left(\frac1n\right)\right]\\
  &=1+\frac{\ln n-1}{n}+o\left(\frac{\ln n}{n}\right)\quad(n\to\infty)
\end{aligned}
\]
得
\[
  \sqrt[n]{a_n}=\left[1+\frac{\ln n-1}{n}+o\left(\frac{\ln n}{n}\right)\right]^n\to+\infty\quad(n\to\infty).
\]
由柯西判别法知 $\sum a_n$ 发散。

\noindent\textbf{例 9.2.8}\quad 试讨论正项级数 $\sum C_n$ 的敛散性, 其中
\[
  C_n=\begin{cases}
    a^{(n+1)/2}, & n\text{ 为奇数},\\
    b^{n/2}, & n\text{ 为偶数}
  \end{cases}\quad(0<a<b<1).
\]

\noindent\textbf{解}\quad 由于
\[
  \overline{\lim}_{n\to\infty}\sqrt[n]{C_n}\leq\lim_{n\to\infty}\sqrt[n]{b^{n/2}}=b^{1/2}<1,
\]
由柯西判别法知级数 $\sum C_n$ 是收敛的。

\noindent\textbf{注}\quad 注意到
\[
  \overline{\lim}_{n\to\infty}\frac{C_{n+1}}{C_n}
  =\lim_{n\to\infty}\frac{b^{(n+1)/2}}{a^{(n+1)/2}}=+\infty,
\]
而
\[
  \underline{\lim}_{n\to\infty}\frac{C_{n+1}}{C_n}
  =\lim_{n\to\infty}\frac{a^{(n+2)/2}}{b^{n/2}}=0,
\]
因此达朗贝尔判别法无法判别 $\sum C_n$ 的敛散性. 值得注意的是, 在某些情况下, 用达朗贝尔判别法则较为简便, 如前面的例 9.2.5。

\subsection{拉贝判别法}

以下我们考虑以 $\sum a_n=\sum\dfrac1{n^p}(p>0)$ 作为比较标准时的判别法. 考查比值 $\dfrac{a_n}{a_{n+1}}$, 读者不难找到规律并得到以下的判别法。

\noindent\textbf{定理 9.2.5 (拉贝(Raabe)判别法)}\quad 设 $\sum a_n$ 为一正项级数。

(1) 若 $\lim\limits_{n\to\infty}n\left(\dfrac{a_n}{a_{n+1}}-1\right)=r>1$, 则 $\sum a_n$ 收敛。

(2) 若 $\overline{\lim}_{n\to\infty}n\left(\dfrac{a_n}{a_{n+1}}-1\right)=r'<1$, 则 $\sum a_n$ 发散。

\noindent\textbf{证明}\quad (1) 取 $r_1(r>r_1>1)$, 则存在 $N_1\in\mathbb{N}$, 当 $n\geq N_1$ 时, 有
\[
  \frac{a_n}{a_{n+1}}>1+\frac{r_1}{n}.
\]
再取 $r_2(r_1>r_2>1)$, 由于 $f(x)=1+r_1x-(1+x)^{r_2}$ 满足 $f(0)=0$, 且 $f'(x)=r_1-r_2(1+x)^{r_2-1}$ 在 $x=0$ 的某个邻域内恒大于 0, 因此存在 $N_2\geq N_1$, 当 $n\geq N_2$ 时, 有
\[
  \frac{a_n}{a_{n+1}}>1+\frac{r_1}{n}>\left(1+\frac1n\right)^{r_2}
  =\frac{(n+1)^{r_2}}{n^{r_2}}.
\]
上式等价于
\[
  (n+1)^{r_2}a_{n+1}<n^{r_2}a_n.
\]
因此当 $n>N_2$ 时, 我们有
\[
  a_n<\frac{N_2^{r_2}a_{N_2}}{n^{r_2}}=\frac{C}{n^{r_2}}\quad(C=N_2^{r_2}a_{N_2}).
\]
由比较判别法, 我们推知 $\sum a_n$ 是收敛的。

(2) $r'<1$, 则存在 $N\in\mathbb{N}$, 当 $n\geq N$ 时, 有
\[
  \frac{a_n}{a_{n+1}}<1+\frac1n=\frac{n+1}{n}.
\]
因此当 $n\geq N$ 时, 有
\[
  na_n<(n+1)a_{n+1},
\]
进一步我们有 $na_n>Na_N$, 即
\[
  a_n>\frac{Na_N}{n}=\frac{C}{n}\quad(C=Na_N).
\]
由比较判别法, 我们推知 $\sum a_n$ 是发散的. 证毕。

\noindent\textbf{例 9.2.9}\quad 讨论正项级数 $\sum\limits_{n=1}^{+\infty}\dfrac{(2n-1)!!}{(2n)!!}$ 的敛散性。

\noindent\textbf{解}\quad 记 $a_n=\dfrac{(2n-1)!!}{(2n)!!}$, 则有
\[
\begin{aligned}
  n\left(\frac{a_n}{a_{n+1}}-1\right)
  &=n\left[\frac{\dfrac{(2n-1)!!}{(2n)!!}}{\dfrac{(2n+1)!!}{(2n+2)!!}}-1\right]\\
  &=n\left(\frac{2n+2}{2n+1}-1\right)\\
  &=\frac{n}{2n+1}\to\frac12.
\end{aligned}
\]
由拉贝判别法知, 级数 $\sum\dfrac{(2n-1)!!}{(2n)!!}$ 发散。

\noindent\textbf{注}\quad 若对例 9.2.9 中的级数用达朗贝尔与柯西判别法, 则无法确定该级数的敛散性。

从定理 9.2.5 的证明中可以看出, 若 (2) 中的条件用以下条件替换其结论也成立: 存在 $N\in\mathbb{N}$, 当 $n>N$ 时, 有
\[
  n\left(\frac{a_n}{a_{n+1}}-1\right)\leq1.
\]

细心的读者不难发现, 以上的一些正项级数判别法实际上可以通过几何级数或级数 $\sum\dfrac1{n^p}$ 中的前后两项的比较找出规律而得到. 按照这种思想, 读者可以从研究其他的``标准级数''来发现一些规律从而找到一些新的判别法。

\subsection{柯西积分判别法}

在许多正项级数 $\sum a_n$ 中, 序列 $\{a_n\}$ 是单调下降趋于零的. $\{a_n\}$ 作为定义在正整数集 $\mathbb{N}$ 上的一个函数, 它的图像是 $Oxy$ 平面内的一个点集. 任取一个 $[1,+\infty)$ 上的单调下降连续函数 $f(x)$, 使得 $f(n)=a_n(n=1,2,\cdots)$(见图 9.2.1)。

我们可以建立如下级数敛散性与无穷积分敛散性的关系。

\noindent\textbf{定理 9.2.6}\quad 设函数 $f(x)$ 在 $[1,+\infty)$ 上单调下降趋于零, 记 $a_n=f(n)(n=1,2,\cdots)$, 则正项级数 $\sum a_n$ 收敛的充分必要条件是无穷积分 $\int_1^{+\infty}f(x)\,dx$ 收敛。

\noindent\textbf{证明}\quad 对任何正整数 $n$, 当 $n\leq x\leq n+1$ 时, 有 $f(n+1)\leq f(x)\leq f(n)$, 因此有
\[
  a_{n+1}=\int_n^{n+1}f(n+1)\,dx\leq\int_n^{n+1}f(x)\,dx\leq\int_n^{n+1}f(n)\,dx=a_n.
\]
对上面不等式关于 $n$ 求和得
\[
  \sum_{n=2}^{+\infty}a_n\leq\int_1^{+\infty}f(x)\,dx
  =\sum_{n=1}^{+\infty}\int_n^{n+1}f(x)\,dx
  \leq\sum_{n=1}^{+\infty}a_n.
\]
因此当 $\int_1^{+\infty}f(x)\,dx$ 收敛时, $\sum\limits_{n=2}^{+\infty}a_n$ 有上界 $\int_1^{+\infty}f(x)\,dx$, 从而 $\sum\limits_{n=1}^{+\infty}a_n$ 收敛. 反之, 当 $\sum a_n$ 收敛时, 无穷积分 $\int_1^{+\infty}f(x)\,dx$ 有上界 $\sum a_n$, 从而积分收敛. 证毕。

在广义积分中, 我们讨论过以下积分的敛散性:
\[
  \int_1^{+\infty}\frac{dx}{x^p},\qquad
  \int_2^{+\infty}\frac{dx}{x^p\ln^q x}.
\]
由定理 9.2.6 知, $\sum\dfrac1{n^p}$ 当 $p>1$ 时收敛, 当 $p\leq1$ 时发散. 对于级数 $\sum\limits_{n=2}^{+\infty}\dfrac1{n^p\ln^q n}$, 当 $p>1$ 时, 对任何 $q$, 该级数都收敛; 当 $p=1$ 时, 对 $q>1$ 该级数收敛, 而对其他情形该级数发散。

值得注意的是, 对于正项级数
\[
  \sum_{n=2}^{+\infty}a_n=\sum_{n=2}^{+\infty}\frac1{n^p\ln^q n},
\]
用柯西判别法及拉贝判别法都无法判别它的敛散性。

\noindent\textbf{例 9.2.10}\quad 设正项级数 $\sum a_n$ 收敛, 记 $r_n=\sum\limits_{k=n}^{+\infty}a_k(n=1,2,\cdots)$, 证明: 当 $p<1$ 时, 正项级数 $\sum\dfrac{a_n}{r_n^p}$ 收敛。

\noindent\textbf{证明}\quad 设 $a_0=0$, 并记 $S=\sum a_n$. 由
\[
  r_n=S-\sum_{k=0}^{n-1}a_k\quad(n=1,2,\cdots)
\]
及 $\lim\limits_{n\to\infty}\sum\limits_{k=0}^{n-1}a_k=S$ 知, $\{r_n\}$ 单调下降趋于零, 因此我们有
\[
  S=r_1\geq r_2\geq r_3\geq r_4\geq\cdots,
\]
且 $r_n>0$。

注意到对于 $\forall n\in\mathbb{N}$, 有
\[
  \frac{a_n}{r_n^p}\leq\int_{r_{n+1}}^{r_n}\frac{dx}{x^p},
\]
从而当 $p<1$ 时, 对于 $\forall n\in\mathbb{N}$, 我们有
\[
  \sum_{k=1}^n\frac{a_k}{r_k^p}
  \leq\sum_{k=1}^n\int_{r_{k+1}}^{r_k}\frac{dx}{x^p}
  <\int_0^S\frac{dx}{x^p}<+\infty.
\]
因此 $\sum\dfrac{a_n}{r_n^p}$ 收敛。

上例告诉我们, 当正项级数 $\sum a_n(a_n>0)$ 收敛时, 我们总能找到另一个收敛的正项级数 $\sum b_n$, 使得 $\lim\limits_{n\to\infty}\dfrac{b_n}{a_n}=+\infty$。

\section{任意项级数}

\subsection{交错级数的敛散性}

对任意项级数 $\sum\limits_{n=1}^{+\infty}a_n$ 而言, 若要讨论它是否绝对收敛时, 我们可以转化为讨论正项级数的敛散性问题. 因此对于任意项级数, 我们现在要讨论的问题是如何判定它们的条件收敛性。

数项级数 $\sum\limits_{n=1}^{+\infty}a_n$ 收敛的一个必要条件是 $\lim\limits_{n\to\infty}a_n=0$. 由于该条件不是充分的, 因此我们还需附加一定的条件才能保证 $\sum\limits_{n=1}^{+\infty}a_n$ 的收敛性. 我们有以下的定理。

\noindent\textbf{定理 9.3.1}\quad 设数项级数 $\sum\limits_{n=1}^{+\infty}a_n$ 满足条件:

(1) $\lim\limits_{n\to\infty}a_n=0$;

(2) 存在 $N\in\mathbb{N}$, 使得在 $\sum\limits_{n=1}^{+\infty}a_n$ 中加上一些括号, 并且在每个括号中的加数个数 $\leq N$, 得到的级数 $\sum\limits_{k=1}^{+\infty}b_k$ 是收敛的,

则级数 $\sum\limits_{n=1}^{+\infty}a_n$ 收敛。

\noindent\textbf{证明}\quad 记 $S_n=\sum\limits_{k=1}^n a_k(n=1,2,\cdots)$. 由于加括号后的级数收敛, 因此存在 $\{S_n\}$ 的一个子列 $\{S_{n_k}\}$ 是收敛的. 对任意的正整数 $n>n_1$, 必存在 $n_k$, 使得 $n_k\leq n<n_{k+1}$, 因此有
\[
  |S_n-S_{n_k}|\leq\sum_{j=1}^N|a_{n_k+j}|.
\]
注意到 $n\to\infty$ 时必有 $n_k\to\infty$, 由 $\lim\limits_{j\to\infty}\sum\limits_{k=j}^{j+N}|a_k|=0$, 我们推出 $\lim\limits_{n\to\infty}S_n$ 存在, 即 $\sum\limits_{n=1}^{+\infty}a_n$ 收敛. 证毕。

下面我们来研究一类特殊的任意项级数——交错级数. 设 $\{a_n\}$ 是一个序列, 且对 $\forall n\in\mathbb{N}$, $a_n\geq0$, 我们称级数 $\sum\limits_{n=1}^{+\infty}(-1)^{n-1}a_n$ 或 $\sum\limits_{n=1}^{+\infty}(-1)^na_n$ 为一个交错级数. 由定理 9.3.1, 可推出下面重要的结论。

\noindent\textbf{推论 (莱布尼茨交错级数判别法)}\quad 设 $\{a_n\}$ 是一个单调序列, 并且 $\lim\limits_{n\to\infty}a_n=0$, 则交错级数 $\sum\limits_{n=1}^{+\infty}(-1)^{n-1}a_n$ 收敛。

\noindent\textbf{证明}\quad 不妨设序列 $\{a_n\}$ 单调下降趋于零, 从而必有 $a_n\geq0(n=1,2,\cdots)$. 显然 $\sum\limits_{k=1}^{+\infty}(a_{2k-1}-a_{2k})$ 是在 $\sum\limits_{n=1}^{+\infty}(-1)^{n-1}a_n$ 中加上括号后得到的正项级数. 记该正项级数的前 $n$ 项部分和为 $S_n$, 我们有
\[
\begin{aligned}
  0\leq S_n
  &=(a_1-a_2)+(a_3-a_4)+\cdots+(a_{2n-1}-a_{2n})\\
  &\leq a_1-(a_2-a_3)-(a_4-a_5)-\cdots-(a_{2n-2}-a_{2n-1})-a_{2n}\\
  &\leq a_1,
\end{aligned}
\]
因此 $\sum\limits_{k=1}^{+\infty}(a_{2k-1}-a_{2k})$ 收敛. 由定理 9.3.1 知 $\sum\limits_{n=1}^{+\infty}(-1)^{n-1}a_n$ 也收敛。

\noindent\textbf{注}\quad 对上述定理的证明略加分析, 我们可以得到以下结论: 设 $\{a_n\}$ 是单调下降趋于零的序列, 则对任意的正整数 $N$, 以下不等式成立:
\[
  0\leq\left|\sum_{n=N}^{+\infty}(-1)^{n-1}a_n\right|\leq a_N.
\]
另外, 我们今后将收敛的交错级数称为莱布尼茨交错级数。

\noindent\textbf{例 9.3.1}\quad 试讨论交错级数 $\sum\limits_{n=2}^{+\infty}\dfrac{(-1)^n}{n^p\ln^q n}(p,q\in\mathbb{R})$ 的敛散性。

\noindent\textbf{解}\quad 从上节知, 当 $p>1$ 时, 对于 $\forall q\in\mathbb{R}$, 该级数绝对收敛; 当 $p=1,q>1$ 时, 该级数也绝对收敛。

当 $0<p<1$ 时, 对于 $\forall q\in\mathbb{R}$, 我们有 $\lim\limits_{n\to\infty}\dfrac1{n^p\ln^q n}=0$. 注意到当 $x$ 充分大时, 有
\[
\begin{aligned}
  (x^{-p}\ln^{-q}x)'
  &=-px^{-p-1}\ln^{-q}x-x^{-p}(-q)(\ln^{-q-1}x)\frac1x\\
  &=x^{-p-1}\ln^{-q}x\left(-p+\frac{q}{\ln x}\right)<0,
\end{aligned}
\]
我们推出, 存在 $N\in\mathbb{N}$, 当 $n>N$ 时, $\left\{\dfrac1{n^p\ln^q n}\right\}$ 单调下降. 因此当 $0<p<1$ 时, 对于 $\forall q\in\mathbb{R}$, $\sum\limits_{n=2}^{+\infty}\dfrac{(-1)^n}{n^p\ln^q n}$ 收敛. 显然此时级数 $\sum\limits_{n=2}^{+\infty}\dfrac1{n^p\ln^q n}$ 发散, 从而原级数条件收敛。

当 $p=0$ 时,
\[
  \sum_{n=2}^{+\infty}\frac{(-1)^n}{n^p\ln^q n}
  =\sum_{n=1}^{+\infty}\frac{(-1)^n}{\ln^q n},
\]
显然该级数当 $q>0$ 时条件收敛, 当 $q\leq0$ 发散。

当 $p<0$ 时, 由于对于 $\forall q\in\mathbb{R}$, 有 $\lim\limits_{n\to\infty}\dfrac1{n^p\ln^q n}=\infty$, 因此 $\sum\limits_{n=2}^{+\infty}\dfrac{(-1)^n}{n^p\ln^q n}$ 发散。

\subsection{狄利克雷判别法和阿贝尔判别法}

如同广义积分一样, 当数项级数具有形式 $\sum\limits_{n=1}^{+\infty}a_nb_n$ 时, 在一些条件下我们也有相应的狄利克雷判别法和阿贝尔判别法。

\noindent\textbf{定理 9.3.2 (狄利克雷判别法)}\quad 设数项级数 $\sum\limits_{n=1}^{+\infty}a_n$ 的部分和序列 $\left\{S_n=\sum\limits_{k=1}^n a_k\right\}$ 是有界的, $\{b_n\}$ 是单调序列且 $\lim\limits_{n\to\infty}b_n=0$, 则数项级数 $\sum\limits_{n=1}^{+\infty}a_nb_n$ 收敛。

\noindent\textbf{证明}\quad 由 $\{S_n\}$ 有界, 存在 $M>0$, 使得
\[
  |S_n|\leq\frac{M}{2}\quad(n=1,2,\cdots),
\]
从而对任意的正整数 $n,p$, 有
\[
  |S_{n+p}-S_n|\leq M.
\]
对于 $\forall\varepsilon>0$, 由 $\lim\limits_{n\to\infty}b_n=0$, 存在 $N\in\mathbb{N}$, 当 $n>N$ 时, 有
\[
  |b_n|<\frac{\varepsilon}{6M}.
\]
因此当 $n>N$ 时, 对任意正整数 $p$, 利用阿贝尔变换(引理 7.5.4)得
\[
\begin{aligned}
  \left|\sum_{k=n+1}^{n+p}a_kb_k\right|
  &=\left|b_{n+p}(S_{n+p}-S_n)-\sum_{k=n+1}^{n+p-1}(S_k-S_n)(b_{k+1}-b_k)\right|\\
  &\leq M\left\{|b_{n+p}|+\sum_{k=n+1}^{n+p-1}|b_{k+1}-b_k|\right\}\\
  &=M\left\{|b_{n+p}|+\left|\sum_{k=n+1}^{n+p-1}(b_{k+1}-b_k)\right|\right\}\\
  &\leq M\{|b_{n+1}|+2|b_{n+p}|\}<\varepsilon.
\end{aligned}
\]
由柯西准则知 $\sum\limits_{n=1}^{+\infty}a_nb_n$ 收敛. 证毕。

\noindent\textbf{定理 9.3.3 (阿贝尔判别法)}\quad 设数项级数 $\sum\limits_{n=1}^{+\infty}a_n$ 收敛, 序列 $\{b_n\}$ 单调有界, 则数项级数 $\sum\limits_{n=1}^{+\infty}a_nb_n$ 收敛。

\noindent\textbf{证明}\quad 由于 $\{b_n\}$ 是单调有界序列, 从而收敛. 设 $\lim\limits_{n\to\infty}b_n=b$, 我们有
\[
  \sum_{n=1}^{+\infty}a_nb_n
  =\sum_{n=1}^{+\infty}(b_n-b)a_n+\sum_{n=1}^{+\infty}ba_n.
\]
显然级数 $\sum\limits_{n=1}^{+\infty}ba_n$ 是收敛的, 而在级数 $\sum\limits_{n=1}^{+\infty}(b_n-b)a_n$ 中令 $b'_n=b_n-b(n=1,2,\cdots)$, 则 $\{b'_n\}$ 及 $\sum\limits_{n=1}^{+\infty}a_n$ 满足狄利克雷判别法(定理 9.3.2)的条件, 因此 $\sum\limits_{n=1}^{+\infty}(b_n-b)a_n$ 收敛, 所以 $\sum\limits_{n=1}^{+\infty}a_nb_n$ 收敛. 证毕。

\noindent\textbf{注}\quad 容易看出来布尼茨交错级数判别法可以由狄利克雷判别法直接推出。

\noindent\textbf{例 9.3.2}\quad 试讨论交错级数 $\sum\limits_{n=1}^{+\infty}\dfrac{(-1)^n}{n^{\alpha+\frac1n}}(\alpha>0)$ 的敛散性。

\noindent\textbf{解}\quad 显然, 当 $\alpha>1$ 时, 该级数绝对收敛. 当 $0<\alpha\leq1$ 时, 对 $n\in\mathbb{N}$, 记 $a_n=\dfrac{(-1)^n}{n^\alpha}$ 和 $b_n=\dfrac1{n^{\frac1n}}$. 由莱布尼茨交错级数判别法知 $\sum\limits_{n=1}^{+\infty}a_n$ 收敛. 又记 $n^{\frac1n}=e^{\frac{\ln n}{n}}$, 则容易证明: 当 $n$ 充分大时, 序列 $\left\{\dfrac{\ln n}{n}\right\}$ 是单调下降的. 因此序列 $\left\{\dfrac1{n^{\frac1n}}\right\}$ 当 $n$ 充分大时单调上升, 且有上界 1. 于是由阿贝尔判别法知 $\sum\limits_{n=1}^{+\infty}\dfrac{(-1)^n}{n^{\alpha+\frac1n}}$ 收敛. 由于
\[
  \lim_{n\to\infty}\frac{\dfrac1{n^{\alpha+\frac1n}}}{\dfrac1{n^\alpha}}=1,
\]
而当 $0<\alpha\leq1$ 时, $\sum\limits_{n=1}^{+\infty}\dfrac1{n^\alpha}$ 发散, 从而 $\sum\limits_{n=1}^{+\infty}\dfrac1{n^{\alpha+\frac1n}}$ 发散, 故此时原级数是条件收敛的。

\noindent\textbf{例 9.3.3}\quad 设序列 $\{a_n\}$ 单调且 $\lim\limits_{n\to\infty}a_n=0$, 证明:

(1) 当 $\sum\limits_{n=1}^{+\infty}a_n$ 收敛时, $\sum\limits_{n=1}^{+\infty}a_n\sin nx$ 与 $\sum\limits_{n=1}^{+\infty}a_n\cos nx$ 对一切 $x\in\mathbb{R}$ 都绝对收敛;

(2) 当 $\sum\limits_{n=1}^{+\infty}a_n$ 发散时, $\sum\limits_{n=1}^{+\infty}a_n\sin nx$ 对一切 $x\neq k\pi(k\in\mathbb{Z})$ 条件收敛, 而 $\sum\limits_{n=1}^{+\infty}a_n\cos nx$ 对一切 $x\neq2k\pi(k\in\mathbb{Z})$ 条件收敛。

\noindent\textbf{证明}\quad (1) 对于 $\forall n\in\mathbb{N}$, 我们不妨设 $a_n>0$, 此时显然有
\[
  |a_n\sin nx|\leq a_n,\qquad |a_n\cos nx|\leq a_n.
\]
因此当 $\sum\limits_{n=1}^{+\infty}a_n$ 收敛时, $\sum\limits_{n=1}^{+\infty}a_n\sin nx$ 与 $\sum\limits_{n=1}^{+\infty}a_n\cos nx$ 对一切 $x\in\mathbb{R}$ 都绝对收敛。

(2) 我们先证两个级数在给定的条件下是收敛的. 由狄利克雷判别法, 只要证在所给条件下, 对于 $\forall n\in\mathbb{N}$,
\[
  S_n=\sum_{k=1}^n\sin kx,\qquad S'_n=\sum_{k=1}^n\cos kx,
\]
有界即可. 事实上, 由
\[
\begin{aligned}
  2\sin\frac{x}{2}\cdot S_n
  &=\sum_{k=1}^n\left[\cos\left(k-\frac12\right)x-\cos\left(k+\frac12\right)x\right]\\
  &=\cos\frac{x}{2}-\cos\left(n+\frac12\right)x,
\end{aligned}
\]
当 $x\neq k\pi(k\in\mathbb{Z})$ 时(当 $x=k\pi$ 时, 该级数的每一项均为 0), 对任意的 $n\in\mathbb{N}$, 有
\[
  |S_n|\leq\frac1{\left|\sin\dfrac{x}{2}\right|}.
\]
而对于 $\forall n\in\mathbb{N}$, 由
\[
\begin{aligned}
  2\sin\frac{x}{2}\cdot S'_n
  &=-\sum_{k=1}^n\left[\sin\left(k-\frac12\right)x-\sin\left(k+\frac12\right)x\right]\\
  &=\sin\left(n+\frac12\right)x-\sin\frac{x}{2},
\end{aligned}
\]
当 $x\neq2k\pi(k\in\mathbb{Z})$ 时, 有
\[
  |S'_n|\leq\frac1{\left|\sin\dfrac{x}{2}\right|}.
\]
因此由狄利克雷判别法知, $\sum\limits_{n=1}^{+\infty}a_n\sin nx$ 对一切 $x\in\mathbb{R}$ 都收敛, 而 $\sum\limits_{n=1}^{+\infty}a_n\cos nx$ 对 $x\neq2k\pi(k\in\mathbb{Z})$ 收敛。

当 $\sum\limits_{n=1}^{+\infty}a_n$ 发散时, 注意到当 $x=k\pi(k\in\mathbb{Z})$ 时, $\sum\limits_{n=1}^{+\infty}a_n\sin nx=\sum\limits_{n=1}^{+\infty}0$, 下面假定 $x\neq k\pi(k\in\mathbb{Z})$, 我们证明级数 $\sum\limits_{n=1}^{+\infty}a_n\sin nx$ 条件收敛. 事实上, 由于对于任意的 $n\in\mathbb{N}$, 有
\[
  |a_n\sin nx|\geq |a_n\sin^2 nx|=\frac{a_n}{2}(1-\cos2nx),
\]
而 $\sum\limits_{n=1}^{+\infty}a_n\cos2nx$ 收敛, $\sum\limits_{n=1}^{+\infty}a_n=+\infty$, 因此 $\sum\limits_{n=1}^{+\infty}|a_n\sin nx|=+\infty$, 从而 $\sum\limits_{n=1}^{+\infty}a_n\sin nx$ 条件收敛。

同理, 当 $x\neq2k\pi(k\in\mathbb{Z})$ 且 $\sum\limits_{n=1}^{+\infty}a_n$ 发散时, $\sum\limits_{n=1}^{+\infty}a_n\cos nx$ 条件收敛。

\section{数项级数的性质}

由于数项级数是由一列数相加而成, 人们自然要问: 对级数而言加法结合律、交换律是否还成立? 当讨论两个数项级数相乘的问题时, 我们还会遇到分配律的问题. 在本节中我们将仔细研究这些问题。

\subsection{结合律}

设 $\sum\limits_{n=1}^{+\infty}a_n$ 是一个数项级数, 对该级数加上一些括号并将每个括号中的加数求和后作为一项将得到一个新的级数 $\sum\limits_{k=1}^{+\infty}b_k$. 关于数项级数的结合律, 我们有以下的结论。

\noindent\textbf{定理 9.4.1}\quad 设 $\sum\limits_{n=1}^{+\infty}a_n$ 是收敛级数, 则在 $\sum\limits_{n=1}^{+\infty}a_n$ 中任意加括号后得到的级数 $\sum\limits_{k=1}^{+\infty}b_k$ 也收敛, 并且 $\sum\limits_{n=1}^{+\infty}a_n=\sum\limits_{k=1}^{+\infty}b_k$。

\noindent\textbf{证明}\quad 对于 $\forall n\in\mathbb{N}$, 设 $S_n$ 是 $\sum\limits_{n=1}^{+\infty}a_n$ 前 $n$ 项部分和; 对于 $\forall k\in\mathbb{N}$, 设 $S'_k$ 是 $\sum\limits_{k=1}^{+\infty}b_k$ 的前 $k$ 项的部分和. 容易看出, $\{S'_k\}$ 是 $\{S_n\}$ 的一个子列, 因此当 $\lim\limits_{n\to\infty}S_n$ 存在时, $\lim\limits_{k\to\infty}S'_k$ 也存在, 并且 $\lim\limits_{k\to\infty}S'_k=\lim\limits_{n\to\infty}S_n$. 证毕。

数项级数 $\sum\limits_{n=1}^{+\infty}(-1)^n$ 发散告诉我们, 级数加上括号后得到的级数收敛并不能推出原来的级数收敛。

对于正项级数 $\sum\limits_{n=1}^{+\infty}a_n$, 我们容易看出, 如果在 $\sum\limits_{n=1}^{+\infty}a_n$ 中加上括号后得到的级数是收敛的, 则 $\sum\limits_{n=1}^{+\infty}a_n$ 也收敛. 对于一般的级数来说, 如果在 $\sum\limits_{n=1}^{+\infty}a_n$ 中加上括号, 并且在每个括号中的项都具有相同的符号, 将这样得到的级数记为 $\sum\limits_{k=1}^{+\infty}b_k$, 则 $\sum\limits_{n=1}^{+\infty}a_n$ 收敛的充分必要条件是 $\sum\limits_{k=1}^{+\infty}b_k$ 收敛, 并且当它们收敛时它们的和也相同(见本章习题)。

\subsection{交换律}

交换律对无穷多个数相加是否还成立呢? 为此我们要讨论数项级数 $\sum\limits_{n=1}^{+\infty}a_n$ 与它的各项重排后所得的级数 $\sum\limits_{n=1}^{+\infty}a'_n$ 的关系. 为了精确描述级数的重排, 我们首先来定义正整数集 $\mathbb{N}$ 的重排: 若函数 $f(n)$ 是 $\mathbb{N}$ 到 $\mathbb{N}$ 的一个一一对应, 则称 $f(n)(n\in\mathbb{N})$ 是 $\mathbb{N}$ 的一个重排。

现设 $\sum\limits_{n=1}^{+\infty}a_n$ 是一个数项级数. 我们称 $\sum\limits_{n=1}^{+\infty}a'_n$ 是它的一个重排, 如果存在 $\mathbb{N}$ 的一个重排 $f(n)$, 使得 $a'_n=a_{f(n)}(n=1,2,\cdots)$。

对于一般的数项级数来说, 我们有以下的结果。

\noindent\textbf{定理 9.4.2}\quad 设 $\sum\limits_{n=1}^{+\infty}a_n$ 为一个数项级数, $f(n):\mathbb{N}\to\mathbb{N}$ 为一个重排, 再设 $\exists M>0$, 使得对于 $\forall n\in\mathbb{N}$, 有 $|f(n)-n|\leq M$, 则级数 $\sum\limits_{n=1}^{+\infty}a_n$ 收敛的充分必要条件是级数 $\sum\limits_{n=1}^{+\infty}a_{f(n)}$ 收敛, 并且当它们收敛时, 有 $\sum\limits_{n=1}^{+\infty}a_n=\sum\limits_{n=1}^{+\infty}a_{f(n)}$。

\noindent\textbf{证明}\quad 必要性\quad 设 $\sum\limits_{n=1}^{+\infty}a_n$ 收敛, 下证 $\sum\limits_{n=1}^{+\infty}a_{f(n)}$ 也收敛. 由于 $\sum\limits_{n=1}^{+\infty}a_n$ 收敛, 我们有 $\lim\limits_{n\to\infty}a_n=0$. 又因对于 $\forall n\in\mathbb{N}$, 有 $|f(n)-n|\leq M$, 从而有
\[
  \lim_{n\to\infty}\sum_{j=-M}^{M}|a_{n+j}|=0.
\]
设 $\{S_n\}$ 和 $\{S'_n\}$ 分别为 $\sum\limits_{n=1}^{+\infty}a_n$ 和 $\sum\limits_{n=1}^{+\infty}a_{f(n)}$ 的部分和序列, 容易看出
\[
  |S_n-S'_n|\leq\sum_{j=-M}^{M}|a_{n+j}|\to0\quad(n\to\infty).
\]
因此 $\lim\limits_{n\to\infty}S'_n=\lim\limits_{n\to\infty}S_n$, 即 $\sum\limits_{n=1}^{+\infty}a_{f(n)}$ 收敛, 且 $\sum\limits_{n=1}^{+\infty}a_n=\sum\limits_{n=1}^{+\infty}a_{f(n)}$。

充分性\quad 设 $\sum\limits_{n=1}^{+\infty}a_{f(n)}$ 收敛, 显然 $\sum\limits_{n=1}^{+\infty}a_n$ 是 $\sum\limits_{n=1}^{+\infty}a_{f(n)}$ 的一个重排, 其重排函数为 $f^{-1}(n):\mathbb{N}\to\mathbb{N}$. 容易看出, 对于 $\forall n\in\mathbb{N}$, 有
\[
  |f^{-1}(f(n))-f(n)|=|n-f(n)|\leq M.
\]
由必要性的证明知 $\sum\limits_{n=1}^{+\infty}a_n$ 收敛, 且 $\sum\limits_{n=1}^{+\infty}a_n=\sum\limits_{n=1}^{+\infty}a_{f(n)}$. 证毕。

\noindent\textbf{注}\quad 定理 9.4.2 告诉我们: 若一个级数交换各项的位置, 只要每项离原来所在的位置不要太远, 则不会改变级数的敛散性。

对于绝对收敛的级数, 我们有下面满意的结果:

\noindent\textbf{定理 9.4.3}\quad 设 $\sum\limits_{n=1}^{+\infty}a_n$ 是绝对收敛的, 则它的任意一个重排 $\sum\limits_{n=1}^{+\infty}a_{f(n)}$ 也绝对收敛, 并且 $\sum\limits_{n=1}^{+\infty}a_n=\sum\limits_{n=1}^{+\infty}a_{f(n)}$。

\noindent\textbf{证明}\quad 先证 $\sum\limits_{n=1}^{+\infty}|a_{f(n)}|$ 收敛并且 $\sum\limits_{n=1}^{+\infty}|a_{f(n)}|=\sum\limits_{n=1}^{+\infty}|a_n|$. 记 $\{S'_n\}$ 为正项级数 $\sum\limits_{n=1}^{+\infty}|a_{f(n)}|$ 的部分和序列, 显然, 对于 $\forall n\in\mathbb{N}$, 有
\[
  S'_n\leq\sum_{n=1}^{+\infty}|a_n|.
\]
这说明了 $\sum\limits_{n=1}^{+\infty}|a_{f(n)}|$ 是收敛的, 即 $\sum\limits_{n=1}^{+\infty}a_{f(n)}$ 是绝对收敛的. 再令 $n\to\infty$, 得
\[
  \sum_{n=1}^{+\infty}|a_{f(n)}|\leq\sum_{n=1}^{+\infty}|a_n|.
\]
由于 $\sum\limits_{n=1}^{+\infty}a_n$ 是 $\sum\limits_{n=1}^{+\infty}a_{f(n)}$ 的一个重排, 我们同样有
\[
  \sum_{n=1}^{+\infty}|a_n|\leq\sum_{n=1}^{+\infty}|a_{f(n)}|.
\]
因此有
\[
  \sum_{n=1}^{+\infty}|a_n|=\sum_{n=1}^{+\infty}|a_{f(n)}|.
\]
为了证明 $\sum\limits_{n=1}^{+\infty}a_n=\sum\limits_{n=1}^{+\infty}a_{f(n)}$, 对于任意的 $n\in\mathbb{N}$, 我们定义
\[
  a_n^+=
  \begin{cases}
    a_n, & a_n>0,\\
    0, & a_n\leq0,
  \end{cases}
  \qquad
  a_n^-=
  \begin{cases}
    -a_n, & a_n<0,\\
    0, & a_n\geq0.
  \end{cases}
  \tag{9.4.1}
\]
显然对于 $\forall n\in\mathbb{N}$, 有 $a_n^+\leq|a_n|$, $a_n^-\leq|a_n|$. 因此 $\sum\limits_{n=1}^{+\infty}a_n^+$ 与 $\sum\limits_{n=1}^{+\infty}a_n^-$ 均为收敛的正项级数, 且有 $\sum\limits_{n=1}^{+\infty}a_n^+=\sum\limits_{n=1}^{+\infty}a_{f(n)}^+$ 和 $\sum\limits_{n=1}^{+\infty}a_n^-=\sum\limits_{n=1}^{+\infty}a_{f(n)}^-$. 由此得
\[
  \sum_{n=1}^{+\infty}a_n
  =\sum_{n=1}^{+\infty}a_n^+-\sum_{n=1}^{+\infty}a_n^-
  =\sum_{n=1}^{+\infty}a_{f(n)}^+-\sum_{n=1}^{+\infty}a_{f(n)}^-
  =\sum_{n=1}^{+\infty}a_{f(n)}.
\]
证毕。

对于条件收敛的级数, 则上述定理可能不成立. 我们来考查以下例子. 我们知道数项级数
\[
  \sum_{n=1}^{+\infty}\frac{(-1)^{n-1}}{n}=1-\frac12+\frac13-\frac14+\cdots
\]
是条件收敛的级数. 利用 $1+\dfrac12+\cdots+\dfrac1n=\ln n+c+\varepsilon_n$, 其中 $c$ 为欧拉常数, $\varepsilon_n\to0(n\to\infty)$, 我们得到级数 $\sum\limits_{n=1}^{+\infty}\dfrac{(-1)^{n-1}}{n}$ 前 $2n$ 项的部分和
\[
\begin{aligned}
  S_{2n}
  &=1-\frac12+\frac13-\cdots-\frac1{2n}\\
  &=\left(1+\frac12+\frac13+\cdots+\frac1{2n}\right)-2\left(\frac12+\frac14+\cdots+\frac1{2n}\right)\\
  &=\ln2n+\varepsilon_{2n}+c-(\ln n+c+\varepsilon_n)\\
  &=\ln2+\varepsilon_{2n}-\varepsilon_n\to\ln2\quad(n\to\infty).
\end{aligned}
\]
再由级数 $\sum\limits_{n=1}^{+\infty}\dfrac{(-1)^{n-1}}{n}$ 前 $2n+1$ 项的部分和
\[
  S_{2n+1}=S_{2n}+\frac1{2n+1}\to\ln2\quad(n\to\infty),
\]
知 $\sum\limits_{n=1}^{+\infty}\dfrac{(-1)^{n-1}}{n}$ 收敛到 $\ln2$。

我们将此级数重排, 使得每两个正项后面跟一个负项, 即
\[
  1+\frac13-\frac12+\frac15+\frac17-\frac14+\cdots+\frac1{4n-3}+\frac1{4n-1}-\frac1{2n}+\cdots.
\]
对于 $n\in\mathbb{N}$, 设它的前 $n$ 项部分和 $S'_n$, 则
\[
\begin{aligned}
  S'_{3n}
  &=\sum_{k=1}^n\left(\frac1{4k-3}+\frac1{4k-1}-\frac1{2k}\right)\\
  &=\sum_{k=1}^{4n}\frac1k-\sum_{k=1}^n\left(\frac1{4k-2}+\frac1{4k}+\frac1{2k}\right).
\end{aligned}
  \tag{9.4.2}
\]
而
\[
  \sum_{k=1}^n\left(\frac1{4k-2}+\frac1{4k}+\frac1{2k}\right)
  =\frac12\sum_{k=1}^{2n}\frac1k+\frac12\sum_{k=1}^n\frac1k,
\]
将它与 $1+\dfrac12+\cdots+\dfrac1n=\ln n+c+\varepsilon_n$ 代入 (9.4.2), 得
\[
  S'_{3n}
  =\ln4n+c+\varepsilon_{4n}
  -\frac12(\ln2n+c+\varepsilon_{2n})
  -\frac12(\ln n+c+\varepsilon_n)\to\frac32\ln2\quad(n\to\infty).
\]
再由
\[
  S'_{3n+1}=S'_{3n}+\frac1{4n+1}\to\frac32\ln2\quad(n\to\infty),
\]
\[
  S'_{3n+2}=S'_{3n}+\frac1{4n+1}+\frac1{4n+3}\to\frac32\ln2\quad(n\to\infty),
\]
我们推知重排后的级数收敛到 $\dfrac32\ln2$. 这说明了对于条件收敛的级数, 重排有可能改变级数的和。

对于条件收敛的级数, 我们有下面的一般性结果。

\noindent\textbf{定理 9.4.4 (黎曼定理)}\quad 设数项级数 $\sum\limits_{n=1}^{+\infty}a_n$ 条件收敛, 则对任意的 $S(S$ 为有限实数或者 $\pm\infty)$, 都存在 $\sum\limits_{n=1}^{+\infty}a_n$ 的重排 $\sum\limits_{n=1}^{+\infty}a_{f(n)}$, 使得
\[
  \sum_{n=1}^{+\infty}a_{f(n)}=S.
\]

\noindent\textbf{证明}\quad 对于条件收敛级数 $\sum\limits_{n=1}^{+\infty}a_n$, 我们有 $\lim\limits_{n\to\infty}a_n=0$ 和
\[
  \sum_{n=1}^{+\infty}a_n^+=+\infty,\qquad \sum_{n=1}^{+\infty}a_n^-=+\infty,
\]
其中 $a_n^+$, $a_n^-(n=1,2,\cdots)$ 由 (9.4.1) 定义。

先设 $0\leq S<+\infty$, 因此存在正整数 $p_1$, 使得
\[
  S<\sum_{k=1}^{p_1}a_k^+\leq S+a_{p_1}^+,
\]
同理存在正整数 $q_1$, 使得
\[
  S-a_{q_1}^-\leq\sum_{k=1}^{p_1}a_k^+-\sum_{k=1}^{q_1}a_k^-<S.
\]
由归纳法, 我们可以得到 $\sum\limits_{n=1}^{+\infty}a_n$ 的一个重排
\[
\begin{aligned}
  &(a_1^++\cdots+a_{p_1}^+)-(a_1^-+\cdots+a_{q_1}^-)\\
  &\quad +(a_{p_1+1}^++\cdots+a_{p_2}^+)-(a_{q_1+1}^-+\cdots+a_{q_2}^-)+\cdots\\
  &\quad +(a_{p_{k-1}+1}^++\cdots+a_{p_k}^+)-(a_{q_{k-1}+1}^-+\cdots+a_{q_k}^-)+\cdots,
\end{aligned}
\]
使得它满足
\[
\begin{aligned}
  S
  &< (a_1^++\cdots+a_{p_1}^+)-(a_1^-+\cdots+a_{q_1}^-)\\
  &\quad +(a_{p_1+1}^++\cdots+a_{p_2}^+)-(a_{q_1+1}^-+\cdots+a_{q_2}^-)\\
  &\quad +\cdots+(a_{p_{k-1}+1}^++\cdots+a_{p_k}^+)\\
  &\leq S+a_{p_k}^+
\end{aligned}
\]
和
\[
\begin{aligned}
  S-a_{q_k}^-
  &\leq (a_1^++\cdots+a_{p_1}^+)-(a_1^-+\cdots+a_{q_1}^-)\\
  &\quad +(a_{p_1+1}^++\cdots+a_{p_2}^+)-(a_{q_1+1}^-+\cdots+a_{q_2}^-)\\
  &\quad +\cdots+(a_{p_{k-1}+1}^++\cdots+a_{p_k}^+)-(a_{q_{k-1}+1}^-+\cdots+a_{q_k}^-)\\
  &<S.
\end{aligned}
\]
由 $q_k\geq k$, $p_k\geq k$, 知
\[
  \lim_{k\to\infty}a_{p_k}^+=\lim_{k\to\infty}a_{q_k}^-=0.
\]
因此去掉括号得到重排后的级数也收敛到 $S$(见本章习题 13)。

当 $S<0$ 时, 由于 $\sum\limits_{n=1}^{+\infty}(-a_n)$ 也条件收敛, 因此存在重排 $\sum\limits_{n=1}^{+\infty}(-a_{f(n)})$ 收敛于 $-S$, 从而 $\sum\limits_{n=1}^{+\infty}a_{f(n)}$ 收敛于 $S$。

当 $S=+\infty$ 时, 存在正整数 $p_1$, 使得
\[
  \sum_{k=1}^{p_1}a_k^+>1+a_1^-,
\]
同样存在正整数 $p_2$, 使得
\[
  \sum_{k=1}^{p_1}a_k^++\sum_{k=p_1+1}^{p_2}a_k^+>2+a_1^-+a_2^-.
\]
继续下去, 由于 $\sum\limits_{n=1}^{+\infty}a_n^+=+\infty$, 我们可取正整数 $p_3,p_4,\cdots,p_k$, 使得
\[
  \sum_{k=1}^{p_1}a_k^++\sum_{k=p_1+1}^{p_2}a_k^++\cdots+\sum_{k=p_{k-1}+1}^{p_k}a_k^+
  > k+a_1^-+a_2^-+\cdots+a_k^-,
\]
因此我们得到级数的重排
\[
  a_1^++\cdots+a_{p_1}^+-a_1^-+a_{p_1+1}^++\cdots+a_{p_2}^+-a_2^-+\cdots
  +a_{p_{k-1}+1}^++\cdots+a_{p_k}^+-a_k^-+\cdots.
\]
不难看出该级数必发散到 $+\infty$。

当 $S=-\infty$ 时, 利用上述结果知, $\sum\limits_{n=1}^{+\infty}(-a_n)$ 存在重排, 使得它发散到 $+\infty$, 因此 $\sum\limits_{n=1}^{+\infty}a_n$ 存在重排, 使得它发散到 $-\infty$. 证毕。

\subsection{级数的乘法（分配律）}

现在我们来讨论级数的乘法问题. 设 $\sum\limits_{n=1}^{+\infty}a_n$ 与 $\sum\limits_{n=1}^{+\infty}b_n$ 为两个数项级数, 类似于有限项求和的乘法规则, 作出两级数各项所有可能的乘积 $a_kb_j(k=1,2,\cdots,j=1,2,\cdots)$, 将它们排成以下的无穷矩阵(称为这两个级数的乘积矩阵):
\[
  \begin{pmatrix}
    a_1b_1 & a_1b_2 & \cdots & a_1b_n & \cdots\\
    a_2b_1 & a_2b_2 & \cdots & a_2b_n & \cdots\\
    \cdots & \cdots & \cdots & \cdots & \cdots\\
    a_nb_1 & a_nb_2 & \cdots & a_nb_n & \cdots\\
    \cdots & \cdots & \cdots & \cdots & \cdots
  \end{pmatrix}
\]
任何将该矩阵中的元素排成一个序列的方法都将得到 $\sum\limits_{n=1}^{+\infty}a_n$ 与 $\sum\limits_{n=1}^{+\infty}b_n$ 乘积的一个表示. 显然, 这个矩阵中元素的排序方式是多种多样的, 我们常用的有对角线及正方形方法。

在用对角线方法排序时, $\sum\limits_{n=1}^{+\infty}a_n$ 与 $\sum\limits_{n=1}^{+\infty}b_n$ 的乘积为 $\sum\limits_{n=1}^{+\infty}c_n$, 其中
\[
  c_n=\sum_{k+j=n+1}a_kb_j=a_1b_n+a_2b_{n-1}+\cdots+a_nb_1.
\]
用对角线方法得到的级数的乘积也称为 $\sum\limits_{n=1}^{+\infty}a_n$ 与 $\sum\limits_{n=1}^{+\infty}b_n$ 的柯西乘积。

在用正方形方法排序时, $\sum\limits_{n=1}^{+\infty}a_n$ 与 $\sum\limits_{n=1}^{+\infty}b_n$ 的乘积为 $\sum\limits_{n=1}^{+\infty}d_n$, 其中
\[
\begin{aligned}
  d_1&=a_1b_1,\\
  d_2&=a_1b_2+a_2b_2+a_2b_1,\\
  &\cdots\cdots\cdots\\
  d_n&=a_1b_n+a_2b_n+\cdots+a_nb_n+a_nb_{n-1}+\cdots+a_nb_1.
\end{aligned}
\]
分别记 $\sum\limits_{n=1}^{+\infty}a_n$, $\sum\limits_{n=1}^{+\infty}b_n$, $\sum\limits_{n=1}^{+\infty}d_n$ 的前 $n$ 项部分和为 $S_n,S'_n,S''_n$, 则我们有 $S''_n=S_nS'_n$. 因此当 $\sum\limits_{n=1}^{+\infty}a_n$ 与 $\sum\limits_{n=1}^{+\infty}b_n$ 收敛时, $\sum\limits_{n=1}^{+\infty}d_n$ 也收敛, 并且 $\sum\limits_{n=1}^{+\infty}d_n=\sum\limits_{n=1}^{+\infty}a_n\cdot\sum\limits_{n=1}^{+\infty}b_n$。

对于两个绝对收敛的级数的乘积, 我们有以下定理。

\noindent\textbf{定理 9.4.5}\quad 设级数 $\sum\limits_{n=1}^{+\infty}a_n$ 与 $\sum\limits_{n=1}^{+\infty}b_n$ 都绝对收敛, 则其乘积矩阵中所有元素的任何一个排列构成的级数也绝对收敛, 并且它的和为
\[
  \sum_{n=1}^{+\infty}a_n\cdot\sum_{n=1}^{+\infty}b_n.
\]

\noindent\textbf{证明}\quad 设 $\sum\limits_{n=1}^{+\infty}a_{k_n}b_{j_n}$ 为 $\sum\limits_{n=1}^{+\infty}a_n$ 与 $\sum\limits_{n=1}^{+\infty}b_n$ 的乘积矩阵中所有元素的一个排列构成的级数, 其前 $n$ 项部分和记为 $S_n(n=1,2,\cdots)$. 对于 $\forall n\in\mathbb{N}$, 记
\[
  M_n=\max_{1\leq i\leq n}\{k_i,j_i\},
\]
则
\[
  \sum_{i=1}^n|a_{k_i}b_{j_i}|
  \leq\sum_{k=1}^{M_n}|a_k|\cdot\sum_{j=1}^{M_n}|b_j|
  \leq\sum_{k=1}^{+\infty}|a_k|\cdot\sum_{j=1}^{+\infty}|b_j|.
\]
因此 $\sum\limits_{n=1}^{+\infty}a_{k_n}b_{j_n}$ 绝对收敛. 由定理 9.4.3 知, 它的任何重排也绝对收敛并且和不变。

由于正方形方法排列得到的乘积 $\sum\limits_{n=1}^{+\infty}d_n$ 收敛到 $\sum\limits_{n=1}^{+\infty}a_n\cdot\sum\limits_{n=1}^{+\infty}b_n$, 而 $\sum\limits_{n=1}^{+\infty}d_n$ 是在乘积矩阵中所有元素的一个排列构成的级数加上一些括号所得, 因此乘积矩阵中所有元素的任何一个排列构成的级数都收敛到 $\sum\limits_{n=1}^{+\infty}a_n\cdot\sum\limits_{n=1}^{+\infty}b_n$, 即 $\sum\limits_{n=1}^{+\infty}a_{k_n}b_{j_n}$ 收敛到 $\sum\limits_{n=1}^{+\infty}a_n\cdot\sum\limits_{n=1}^{+\infty}b_n$. 证毕。

\noindent\textbf{例 9.4.1}\quad 设级数 $\sum\limits_{n=1}^{+\infty}a_n=\sum\limits_{n=1}^{+\infty}\dfrac{(-1)^{n+1}}{n^\alpha}$, 证明: 当 $\alpha\leq\dfrac12$ 时, 级数 $\sum\limits_{n=1}^{+\infty}a_n$ 与自身的柯西乘积发散。

\noindent\textbf{证明}\quad 设 $\sum\limits_{n=1}^{+\infty}a_n$ 与自身的柯西乘积为 $\sum\limits_{n=1}^{+\infty}c_n$, 则
\[
  c_n=\sum_{k=1}^n\frac{(-1)^{k+n+1-k}}{k^\alpha(n+1-k)^\alpha}
  =(-1)^{n+1}\sum_{k=1}^n\frac1{k^\alpha(n+1-k)^\alpha}.
\]
由于
\[
  \frac1{k^\alpha(n+1-k)^\alpha}
  \geq\frac1{k^{\frac12}(n+1-k)^{\frac12}}
  \geq\frac2{n+1},
\]
因此对于 $\forall n\in\mathbb{N}$, 有
\[
  |c_n|\geq\frac{2n}{n+1}\geq1\quad(n=1,2,\cdots).
\]
所以 $\sum\limits_{n=1}^{+\infty}c_n$ 发散。

值得注意的是, 该级数自乘的正方形方法排序得到的级数是收敛的. 请读者自证: $\sum\limits_{n=1}^{+\infty}\dfrac{(-1)^{n+1}}{n}$ 与自身的柯西乘积是收敛的。

\section{无穷乘积}

在本节中我们讨论无穷个数相乘的问题. 设 $\{a_n\}$ 是一个序列, 我们用 $\prod\limits_{n=1}^{+\infty}a_n$ 表示该序列中所有项的乘积, 并称它为该序列的无穷乘积. 对于 $\forall n\in\mathbb{N}$, 记
\[
  T_n=\prod_{k=1}^n a_k,
\]
我们称它为该无穷乘积的前 $n$ 项部分积, 而称 $\{T_n\}$ 为该无穷乘积的部分积序列。

对于无穷乘积, 我们容易看出它的一个简单性质: 序列 $\{a_n\}$ 中若有一个元素为零, 则对所有充分大的 $n$ 有 $T_n=0$. 此时部分积序列不能给我们提供任何信息. 在 $\lim\limits_{n\to\infty}T_n=0$ 的情形下, 部分积序列也对无穷乘积的研究很难提供有用的信息. 因此我们给出以下的定义。

\noindent\textbf{定义 9.5.1}\quad 设 $\{a_n\}$ 是一个序列, $\{T_n\}$ 为无穷乘积 $\prod\limits_{n=1}^{+\infty}a_n$ 的部分积序列. 若 $\lim\limits_{n\to\infty}T_n=a(a$ 为非零常数), 则称无穷乘积 $\prod\limits_{n=1}^{+\infty}a_n$ 收敛, 并记 $a=\prod\limits_{n=1}^{+\infty}a_n$. 此时也称 $a$ 为该无穷乘积的积. 若 $\{T_n\}$ 不收敛到一个非零常数, 则称无穷乘积 $\prod\limits_{n=1}^{+\infty}a_n$ 发散。

由上述定义, 我们容易推出下面的定理。

\noindent\textbf{定理 9.5.1}\quad 若无穷乘积 $\prod\limits_{n=1}^{+\infty}a_n$ 收敛, 则 $\lim\limits_{n\to\infty}a_n=1$。

\noindent\textbf{证明}\quad 对于 $n=1,2,\cdots$, 记 $T_n=\prod\limits_{k=1}^n a_k$. 由 $\prod\limits_{n=1}^{+\infty}a_n$ 收敛, 设其积为 $a$, 则
\[
  \lim_{n\to\infty}a_n=\lim_{n\to\infty}\frac{T_n}{T_{n-1}}=\frac{a}{a}=1.
\]
证毕。

根据定理 9.5.1, 若无穷乘积 $\prod\limits_{n=1}^{+\infty}a_n$ 收敛, 则对于 $\forall\varepsilon>0$, 当 $n$ 充分大时, 有
\[
  1-\varepsilon<a_n<1+\varepsilon.
\]
由于改变该乘积前面的有限多项(不能变为 0)不改变无穷乘积的收敛性. 因此收敛的无穷乘积的理论中基本上是对正数序列的无穷乘积的研究。

由定理 9.5.1, 我们可将无穷乘积表成 $\prod\limits_{n=1}^{+\infty}(1+a_n)$ 的形式, 不失一般性, 我们还可以假定 $|a_n|<1(n=1,2,\cdots)$. 在上述假定下, 我们容易建立以下定理。

\noindent\textbf{定理 9.5.2}\quad 无穷乘积 $\prod\limits_{n=1}^{+\infty}(1+a_n)$ 收敛的充分必要条件是数项级数 $\sum\limits_{n=1}^{+\infty}\ln(1+a_n)$ 收敛。

\noindent\textbf{证明}\quad 对于 $\forall n\in\mathbb{N}$, 记 $T_n=\prod\limits_{k=1}^n(1+a_k)$, $S_n=\sum\limits_{k=1}^n\ln(1+a_k)$, 则有 $T_n=e^{S_n}(n=1,2,\cdots)$. 因此由 $e^x$ 的连续性, 当 $\lim\limits_{n\to\infty}S_n=S$ 时, 有 $\lim\limits_{n\to\infty}T_n=e^S$; 反之, 当 $\lim\limits_{n\to\infty}T_n=T>0$ 时, 有 $\lim\limits_{n\to\infty}S_n=\ln T$. 证毕。

\noindent\textbf{推论}\quad 若 $a_n>0(n=1,2,\cdots)$, 则无穷乘积 $\prod\limits_{n=1}^{+\infty}(1+a_n)$ 收敛的充分必要条件是数项级数 $\sum\limits_{n=1}^{+\infty}a_n$ 收敛。

\noindent\textbf{证明}\quad 对于正项级数 $\sum\limits_{n=1}^{+\infty}a_n$, 我们容易看出它与 $\sum\limits_{n=1}^{+\infty}\ln(1+a_n)$ 同时收敛与发散. 由定理 9.5.2 立得推论. 证毕。

\noindent\textbf{例 9.5.1}\quad 讨论无穷乘积 $\prod\limits_{n=1}^{+\infty}\left(1+\dfrac1{n^x}\right)$ 的敛散性。

\noindent\textbf{解}\quad 当 $x\leq1$ 时, 由 $\sum\limits_{n=1}^{+\infty}\dfrac1{n^x}$ 发散, 知 $\prod\limits_{n=1}^{+\infty}\left(1+\dfrac1{n^x}\right)$ 发散. 而当 $x>1$ 时, $\sum\limits_{n=1}^{+\infty}\dfrac1{n^x}$ 收敛, 所以 $\prod\limits_{n=1}^{+\infty}\left(1+\dfrac1{n^x}\right)$ 也收敛. 因此, $\prod\limits_{n=1}^{+\infty}\left(1+\dfrac1{n^x}\right)$ 当 $x>1$ 时收敛, 当 $x\leq1$ 时发散。

\noindent\textbf{例 9.5.2}\quad 求无穷乘积 $\prod\limits_{n=1}^{+\infty}\left[1-\dfrac1{(2n)^2}\right]$ 的积。

\noindent\textbf{解}\quad 对于 $\forall n\in\mathbb{N}$, 无穷乘积 $\prod\limits_{n=1}^{+\infty}\left[1-\dfrac1{(2n)^2}\right]$ 的前 $n$ 项部分积为
\[
  T_n=\prod_{k=1}^n\left[1-\frac1{(2k)^2}\right]
  =\frac{[(2n-1)!!]^2}{[(2n)!!]^2}(2n+1).
\]
在定积分中我们利用 $\sin^n x$ 在 $\left[0,\dfrac{\pi}{2}\right]$ 的积分求出了 $\lim\limits_{n\to\infty}T_n=\dfrac2\pi$, 因此
\[
  \prod_{n=1}^{+\infty}\left[1-\frac1{(2n)^2}\right]=\frac2\pi.
\]
最后我们举一个有趣的例子结束本章。

我们已经知道, 当 $s>1$ 时, 正项级数 $\sum\limits_{n=1}^{+\infty}\dfrac1{n^s}$ 是收敛的. 现在我们来建立该级数与无穷乘积的联系. 对每个固定的 $s>1$, 令
\[
  \zeta(s)=\sum_{n=1}^{+\infty}\frac1{n^s}.
\]
首先我们观察
\[
  \left(1-\frac1{2^s}\right)\zeta(s)
  =\sum_{n=1}^{+\infty}\frac1{n^s}-\sum_{n=1}^{+\infty}\frac1{(2n)^s}
  =\sum\frac1{m^s},
\]
在上式右边的求和中 $m$ 将取遍所有奇数, 即所有不是 2 的整数倍的正整数. 然后再考虑
\[
  \left(1-\frac1{2^s}\right)\left(1-\frac1{3^s}\right)\zeta(s)
  =\sum\frac1{m^s}-\sum\frac1{(3m)^s}
  =\sum\frac1{k^s},
\]
在上式右边求和中 $k$ 取遍所有不是 2,3 的整数倍的正整数. 依此类推, 取所有素数 $2,3,5,7,\cdots,p_n,\cdots$, 则有
\[
  \lim_{n\to\infty}\prod_{k=1}^n\left(1-\frac1{p_k^s}\right)\zeta(s)=1,
\]
即
\[
  \prod_{n=1}^{+\infty}\left(1-\frac1{p_n^s}\right)=\frac1{\zeta(s)}.
\]
由上述等式我们易知素数的个数必定是无穷的. 事实上, 倘若结论不真, 设 $p_N$ 为最大的素数, 则对任意的 $s>1$, 成立
\[
  \prod_{n=1}^N\left(1-\frac1{p_n^s}\right)=\frac1{\zeta(s)}.
\]
上述等式的左边是有限个数相乘, 易见当 $s\to1+0$ 时, 有
\[
  \prod_{n=1}^N\left(1-\frac1{p_n^s}\right)=\prod_{n=1}^N\left(1-\frac1{p_n}\right)>0.
\]
由于 $\lim\limits_{s\to1+0}\zeta(s)=+\infty$, 从而 $\lim\limits_{s\to1+0}\dfrac1{\zeta(s)}=0$. 此矛盾便证明了我们的结论。

\end{document}
